\begin{helper}
For fixed split-outside-blocker data \(G,r,X,\phi,\beta\), and a fixed finite
outside-blocker list \(z_0,\ldots,z_{n-1}\) with
\(z_i\notin X\cup\{r\}\), define the candidate sets by
\[
  B_i=\{x\in X:xz_i\in E(G)\},
\]
There is a genuine mixed-stage event surface on one paired source--target
state space.  For every stage \(j\), let
\[
  R_j=\{i:i<j\},\qquad T_j=\{i:j\le i<n\}.
\]
The stage \(H_j(\omega)\) consists of paired states in which one candidate
\(x\in X\) is active on the source side and its image \(\phi(x)\) is active
on the target side, \(x\) beats all active embedded neighbors in
\(X\cup\{r\}\) on the source side and \(\phi(x)\) beats the corresponding
embedded active neighbors on the target side, every retained common blocker
\(z_i\) with \(i\in T_j\) fails to kill \(x\), and every already-private
blocker \(\beta(x,z_i)\) with \(i\in R_j\) fails to kill \(\phi(x)\).

For a step \(k<n\), there is also a shared unchanged context
\(U_k(\omega)\) on the same paired state space.  It contains the same
single candidate \(x\in B_k\), candidate activation, and
embedded-neighbor winning conditions, the
already-private non-killing tests for exactly \(i<k\), and the retained
common non-killing tests for exactly \(k<i<n\).  It omits only the current
blocker \(i=k\).  The two current boundary events are then obtained from this
same single-witness context by adding either the common current test for
\(z_k\) or the private current test for \(\beta(x,z_k)\).
\end{helper}
\begin{proof}
Define \(R_j\) and \(T_j\) inside the finite type of indices by
\[
  R_j=\{i:i<j\},\qquad T_j=\{i:j\le i\}.
\]
These definitions immediately give the two displayed membership
equivalences.

Define \(H_j(\omega)\) to be exactly the paired-state event displayed in the
statement.  A state of \(H_j(\omega)\) contains a candidate \(x\in X\) such
that the source component has \(x\) active, the target component has
\(\phi(x)\) active, both components satisfy the embedded-neighbor priority
comparisons, the source component satisfies all retained common blocker tests
with indices in \(T_j\), and the target component satisfies all already
private blocker tests with indices in \(R_j\).  Since this is the definition
of \(H_j\), the asserted formula for \(H_j(\omega)\) follows directly.

For \(k<n\), define \(U_k(\omega)\) by choosing one candidate
\(x\in B_k\), by requiring the activation and embedded-neighbor conditions for
that same \(x\), by imposing the private tests for all indices \(i<k\), and
by imposing the common retained tests for all indices \(i>k\).  The equality
\(B_i=\{x\in X:xz_i\in E(G)\}\) is exactly why the private terms
\(\beta(x,z_i)\) are well-defined: \(x\in X\), the displayed adjacency gives
\(xz_i\in E(G)\), and the outside-blocker hypothesis gives
\(z_i\notin X\cup\{r\}\).  This definition uses the same paired state and the
same candidate witness on both sides, so the noncurrent context is literally
shared; it is not a comparison between a source-only cylinder and a
target-only cylinder.  Finally define the source current boundary by adding,
for that same \(x\), the non-killing test for the single common blocker
\(z_k\), and define the target current boundary by adding, for that same
\(x\), the non-killing test for the single private blocker
\(\beta(x,z_k)\).  These are definitions, so the three displayed formulas for
the shared context and the two current boundary events follow immediately.
\end{proof}
